\documentclass[USenglish,twocolumn]{article}

\usepackage[utf8]{inputenc}%(only for the pdftex engine)
%\RequirePackage[no-math]{fontspec}%(only for the luatex or the xetex engine)
\usepackage[big]{dgruyter}
%\graphicspath{{figures/}} % path to folder with figures
\usepackage{xcolor}
\usepackage{float}
\usepackage{lipsum} 
%\usepackage{blindtext}
%\usepackage[showframe]{geometry}
\usepackage{graphicx}
\usepackage[super]{nth}
\usepackage{gensymb}
\usepackage{array}
\usepackage{tabularx}
%\usepackage[usenames, dvipsnames]{color}

\begin{document}
 
\articletype{Research Article{\hfill}Open Access}
\author[1]{Per Lindh}
\author*[2]{Polina Lemenkova}
  \affil[1]{$^1$ Swedish Transport Administration, Gibraltargatan 7, Malmö, Sweden \\ $^2$ Lund University, Division of Building Materials, Box 118, SE- 221-00, Lund, Sweden, E-mail: per.a.lindh@gmail.com}
  \affil[2]{$^3$ Université Libre de Bruxelles (ULB), École polytechnique de Bruxelles (Brussels Faculty of Engineering), Laboratory of Image Synthesis and Analysis (LISA). Campus de Solbosch -- CP 165/57, Av. F. D. Roosevelt 50, B-1050 Brussels, Belgium, E-mail: polina.lemenkova@ulb.be}

  \title{\huge Shear bond and compressive strength of clay stabilized with lime/cement jet grouting and deep mixing: a case of Norvik, Nynäshamn}
  
 \runningtitle{Shear bond and compressive strength of clay stabilized with jet grouting}
 \runningauthor{P. Lindh, P. Lemenkova}

\begin{abstract}
{The tests measured the compression wave velocity (P-wave) and compressive strength (UCS) of soil specimens. Three different binders have been used: 1) Cem II/A, a Portland limestone cement (Portland clinker and limestone), 2) burnt lime; 3) lime kiln dust (LKD). The experimental setup was performed as a three-factor Simplex lattice with a limitation of the LKD content up to 50\%. The mixing amount of binders corresponded to 80, 100 and 120 kg per $m^3$ of clay. For jet grouting, cement (CEM II/A-M), LKD and cement kiln dust (CKD) were used as binders.
\textcolor{red}{\lipsum[1-1]}}
\end{abstract}
  \keywords{keyword; keyword; keyword}
 
%  \classification[PACS]{}
 % \communicated{...}
 % \dedication{...}

  \journalname{Nonlinear Engineering}

\DOI{DOI}
  \startpage{1}
  \received{Aug 31, 2022}
  \revised{..}
  \accepted{..}

  \journalyear{2022}
  \journalvolume{x}
  \journalissue{x}
  
\maketitle

\section{Introduction}
The mechanical performance of soil under load stress has a complex and non-linear behaviour. This is largely controlled by a variety of inner and external factors. Apart from the environmental factors (temperature, humidity), inner properties of soil vary owing to the complex structure of soil as a porous media with varied mineral content, density, viscosity, elasticity, compactness, grain size \cite{Kaellenetal2014}. As the behaviour of weak soil presents a serious problem in  industry, it should be stabilized prior to constructions. 

Analysis of the performance of stabilized soil is a challenging task, since the addition of various combinations of binders may result in different response from soil, depending on its structure and mineral content. For instance, fine- or course-grained will react on various binders differently \cite{Lindh2001}. Many different solutions exist in the literature proposing stabilization methods, each based on different types of soil and binders \cite{Lindh2004,Glendinning,Lindh202201,NIGITHA2022}. It has been shown that universal methods for soil stabilization do not exist and require laboratory-based testing and critical evaluating of the stabilization results.

Testing shear bond and compressive strength of soil aims to estimate the parameters of soil performance under load regarding the content of binders and type of soil. This is primarily intended at evaluation of soil cementation upon stabilization with various binders. This can be done either by pure binders of by blended mixtures with various components of stabilizing agents \citep{HOV2022101162,Bakaiyang}. The jet propagation in the soil is performed from the injection nozzle and results in cementation of soil due to the interaction of particles between the jet of binders and soil. The injected fluid penetrates into the soil, increases pore pressures and reduces grain contact in a soil skeleton \cite{Modonietal2006,Croce}. This results in the increased stiffness, density, strength and improved bearing capacity of soil stabilized by jet grouting \cite{Modonietal2019}.

This paper focuses on the estimation of shear bond and compressive strength of clay collected in the Norvik, Nynäshamn Municipality, located in the central region of Sweden, Stockholm County. Soil of the cold regions is subject to harsh climate and environmental setting and has specific properties caused by geological and climate factors \cite{LindhWinter,Lemenkov2021b,Lemenkov2021a}. Such soil should be stabilized before the use in industrial works. The tested methods included  stabilization of specimens with lime/cement mixtures by jet grouting and deep mixing. Jet grouting is useful for various applications in civil engineering 

\textcolor{red}{\lipsum[1-7]}

\newpage
\section{Study objectives}
The Swedish Geotechnical Institute (SGI) has performed laboratory determinations and evaluation of strength properties on stabilized clay from Norvik, Nynäshamn municipality. The works have been carried out on behalf of the Geomind KB, SMA and CEMEX. The research aim was to test different binder combinations for depth stabilization and for jet injection of clay. The tests included two types: 1) mixing tests for depth stabilization; 2) mixing tests for jet grouting. 

\section{Methodology}

\subsection{Mixture of lime/cement in a binder}
To evaluate binder combinations for deep stabilization, an experimental setup of simplex lattice was applied using existing methodology \citep{Montgomery2019,Myres}. The purpose of using this type of statistical trial planning is to minimize the number of trials while maintaining statistical significance and to ensure that both positive and negative interactions between the binder components are detected. Three different binders have been used in this study: 1) cement type Cem II/A, which is a Portland limestone cement consisting of Portland clinker (K) (content 80\%-94\%) and limestone (LL), 2) burnt lime named Ql in the figures ('quicklime'); 3) lime kiln dust (LKD). 

The test surface in a three-factor simplex lattice is structured as a triangle, Figure \ref{Fig:01}. Pure binders (100\%) are represented as a corners in the triangle and a mixture of two-three binders along the triangle's stripes, as a combination of all the binders inside the triangle of the model. The selected trial points are marked as blue circles, which correspond to tested binder combinations. The mutual ratio between the binders varies depending on where the test point is located within the triangle model. 

In such a way, nine different binder combinations were formed from the mix of these three binders.The full circles represent both test point and extra reference samples, Figure \ref{Fig:01} a). The design according to a restricted simplex lattice is shown in Figure \ref{Fig:01} a) where the restriction consists in maximizing the share of LKD to 50\%. The test was used for the different amounts of binder: 80, 100 and 120 kg per $m^3$ of clay. The experimental design for the amount of binder 100 kg of binder per $m^3$ of clay was supplemented with internal points. The internal points increase the resolution of the response surface. Red ring shows point that were excluded from the analysis, Figure  \ref{Fig:01} b). Each trial point was performed as a double trial with the corners of the triangle representing to 100\% of only one binder. The center of the triangle corresponds to 33\% of the three different binders. The points on the lower border of the triangle correspond to either a pure binder, (cf. Cem II/A or Ql) or, alternatively, to a mixture of 50\% Cem II/A and 50\% Ql. 

\begin{figure}[H]
 \centering
    \includegraphics[width=0.45\textwidth]{Fig_01.jpg}\\
  \caption{Experimental design (a): according to a restricted simplex lattice; (b): supplemented with internal points}\label{Fig:01}
\end{figure}

The curing time for the reference samples (cement/lime 50/50) were determined to be 7, 14, 21 and 28 days. Extra reference tests have been carried out for calibration of seismic measurement in relation to compressive strength on days 7 and 14, Figure \ref{Fig:01} b). Moreover, some extra specimens were produced for compression tests on day 80. The choice of 80 days is explained by the fact that the samples were stored at T=7\degree C. At this storage temperature and a curing time of 80 days, this corresponds to the equivalent of 560 \degree days, which corresponds to 28 days of storage at 20°C. The tests were carried out according to the laboratory guidance of the SGI applied in previous experiments \cite{Lindh2021a} and existing methods of using lime/cement columns for deep column stabilization \cite{CARLSTEN1996355}.

\subsection{Seismic measurements}
Seismic measurements were performed as a free-free resonant column test measuring P-wave velocity of the specimen using existing methodology applied in previous studies \citep{Lindh202210,Lindh202102}. The measurements were used on all specimens modelled in previous step in a simplex test. By producing extra specimens that were used as reference specimens, i.e., for the determination of both P-wave velocity and shear strength on days 7 and 14, a reference to the P-wave velocity of the remaining specimens was used for a correlation to the shear strength of the samples. 

The correlation between the P-waves and UCS has previously been used for surface stabilized specimens \cite{Rydenetal2006,Rydenetal2016}. The P-wave is the axial wave passing through the cylindrical specimen. The vibration was excited by a hammer which stroke on one end of the specimen and measured as P-waves by the accelerator device on another end, respectively. The P-wave has been correlated against the shear strength of the stabilized material. Coupling between the P-wave and the E-module of the material ($E_0$, small strain levels) is represented by the Equation \ref{eq1}.

\begin{equation} \label{eq1}
	E_0=\rho \nu{_p2}=\rho(2L fL)2
\end{equation}

The specifics of these test was that the tested sleeves were not completely filled with soil material. The empty sleeves had eigenfrequencies close to the eigenfrequencies of the test bodies. This meant that the frequency in generated oscillation modes deviated significantly from the natural frequency of the samples located in sleeves. Since the natural frequency of the sleeves dominated, it was these frequencies that were measured, cf. section \emph{4.3 Simulation of seismic measurements}. This meant that none of the measurements on the test specimens in sleeves could be evaluated. In contrast, tests performed on the deformed and trimmed specimens demonstrated robust results.

\subsection{Binder recipe for jet grouting}
Testing of the binder recipes has been carried out partly by the SGI and partly by the CEMEX. The surveys included the three types of tests: 1) Uniaxial Compressive Strength (UCS), 2) Water requirements; 3) Binding time for the four types of soil-binder mixtures: 1) Cem II/A (100\%); 2) Cem II/A/CKD (50\%/50\%); 3) Cem II/A/LKD (50\%/50\%); 4) Cem II/A/LKD/LSfiller (50\%/30\%/20\%). The binder combinations were tested for the jet pillar grouting for reinforcement of soil, and summarised in Table \ref{tab:1} where J=performed test, N=not performed test.

\begin{table}[h!]
\centering
\caption{Tested binder combinations for jet pillar reinforcement. J=performed, N=not performed.}
\label{tab:1}
\setlength\tabcolsep{1pt} 
\smallskip 
\begin{tabular*}{\columnwidth}{@{\extracolsep{\fill}}| l | c | c | c | c |}
 \hline
 Name & A & B & C & D \\ \hline
 Cement (\%) & 100 & 70 & 80 & 70 \\ 
  LKD (\%) & -- & 30 & -- & -- \\ 
 CKD (\%) & -- & -- & 20 & 30 \\ 
 vbt & 1.0 & 1.0 & 1.2 & 1.2 \\
 Clay (17\%) & J & J & J & N \\
 Clay 30\% & J & J & J & N \\
 Clay 37.5\% & J & J & J & N\\
 Clay 58\% & J & J & J & N \\
\bottomrule
\end{tabular*}
\end{table}

The tests carried out by CEMEX comply with the Swedish standards of soil testing EN 196-1 \cite{SIS2016M} and EN 196-3 \cite{SIS2016N}. The conditions involve, among other parameters, that some binders have been mixed with the three parts of the standard sand and half a part water (vbt = 0.5). In the tests carried out at SGI, only water and binder have been mixed and in parts of the tests clay has also been mixed in. This leads to a significant differences in the compressive strength. Furthermore, the test performed by CEMEX is done on prisms and the test performed by SGI is done on cylinders with a diameter of 50 mm and a height of 100 mm.

The binder and water were mixed by the industrial mixer for 5 minutes until a homogeneous mixture was obtained. The samples were then cast into the test cases. In cases where clay was added, the mixing time was increased to 20 minutes in order to obtain a homogeneous mixture. After one to three days of curing, the specimens were de-shaped and trimmed to the correct length. Due to the water separation, the specimens were trimmed mostly on the upper part, because this part was weaker and could be processed more conveniently. We did not detect water separation in the samples containing clay. The testing performed by CEMEX and included in this report is shown in Table \ref{tab:1}.

For jet grouting, cement (CEM II/A-M), LKD and cement kiln dust (CKD) were used as binders and evaluated in terms of their performance. A water binder number (vbt) of 1.0 was used for binder in recipes A and B, and a vbt of 1.2 was used for recipes C and D, respectively. The experiments were carried out with the mixing of clay for the tested binder combinations. The total clay content was 17, 30, 37.5 and 58\% by weight. For binder recipe D, no clay admixtures were carried out due to difficulties in obtaining a homogeneous admixture of clay in that sample and a low compressive strength without clay admixture. Tested binder ratios are reported in Table \ref{tab:1}.

Unlike the measurements performed on stabilized soil for lime/cement column reinforcement, all seismic measurements were performed on the specimens without pipes. Furthermore, temperature measurements were carried out according to the thermos method on Cem II (100\%) with a vct of 1.0 and Cem II/LKD (50/50\%) with a vbt of 1.0. Steel thermoses with a liquid volume of 0.5 l were used for the experiment.

The viscosity was determined in recipes A to D without added clay according to the Marsh cone method. The Marsh cone is a workability test for quality control of cement grouts which evaluates the fluidity based on the plastic viscosity and yield stress from the cone geometry of tested specimens \cite{ROUSSEL2005823}. It considers the rheological parameters and the time of flow in performance of soil-cement pastes, since plastic viscosity vary with composition of binders \cite{Banfill,FERRARIS2001245}. In this case, the mixed clay recipes did not passed the funnel in the Marsh cone. 

%----------------------- checked 27.08. % 

\section{Results}

\subsection{Results from the lime/cement column stabilization}
\subsubsection{Strength determination}
The report shows good interaction effects between lime, cement and LKD for deep stabilization purposes. Most of the binder combinations met the requirements for a minimum shear strength of 200 kPa after 28 days of storage at 7\degree. Furthermore, tests carried out show that two out of four tested recipes for jet grouting passed the requirement of 3 MPa after 28 days of storage even with a clay mixture of 58\%. In the testing of binder combinations, a reference recipe of Cem II/A-V and quicklime was used in a mutual ratio of 50/50. Evaluated shear strength for the reference samples at different curing times is reported in Figure \ref{Fig:02}.

\begin{figure}[H]
 \centering
    \includegraphics[width=0.45\textwidth]{Fig_02.jpg}\\
  \caption{Shear strength as a function of storage time for specimens stabilized with a binder combination of CEM II/A and quicklime (50/50). The amounts of binder corresponded to 80, 100 and 120 kg of binder per cubic meter of clay. The samples are stored in climate rooms with a temperature of 7°C.}\label{Fig:02}
\end{figure}

In order to evaluate the effects of different binder combinations on the shear strength of the stabilized material, a test set-up according to Simplex Lattice was used. The result is presented as a response surface. For a binder quantity corresponding to 80 kg/m3, the result is reported in Figure \ref{Fig:03}. The result shows a clear positive interaction between the various components (Cement, quicklime and LKD) in a binder combination. Positive interaction means that the components generate a higher strength together than would be expected from a linear regression between the effects of the single-slip binders. The model has an explanation rate of 95\%, i.e. the model explains 95\% of the variation in shear strength.

\begin{figure}[H]
 \centering
    \includegraphics[width=0.45\textwidth]{Fig_03.jpg}\\
  \caption{Response surface showing shear strength after 28 days of curing at 7\degree C with varied amount of binder. (a):  80 kg/$m^3$; (b): $100 kg/m^3$; (c): $120 kg/m^3$}\label{Fig:03}
\end{figure}

The test set-up for the amount of binder $100 kg/m^3$ clay contains an extended number of tests. For comparison with the other amounts of binder, an analysis was also performed with the same number of test points as for 80 and 120 kg of binder, respectively, Figure \ref{Fig:03}. The result from this analysis shows a slightly lower positive interaction between the binders. For the mixing amount corresponding to 120 kg of binder per $m^3$ of clay, the response surface was obtained according to Figure \ref{Fig:03}. 

However, this must be compared with the result from the analysis of the extended trial reported in Figure \ref{Fig:04} a). In the trial set-up with more internal trial points, the resolution increases and the model's degree of explanation increases from 95\% to 97.5\%1. The extended model shows a greater interaction between the different binder components. A three-dimensional presentation of the response surface is shown in Figure \ref{Fig:04} b).

\begin{figure}[H]
 \centering
    \includegraphics[width=0.45\textwidth]{Fig_04.jpg}\\
  \caption{Response surface with internal points showing shear strength after 28 days of curing at 7\degree C. The amount of binder $100 kg/m^3$. (a): 2D view; (b): 3D view}\label{Fig:04}
\end{figure}

At this mixing amount, the positive interaction of the binder components was small. The degree of explanation of the response surface was 98\%.

%\textcolor{red}{\lipsum[1-2]}

\subsubsection{Seismic measurements}

The seismic measurements were performed as P-wave measurements according to the free-free resonant column method. For the dimensioning of K/C columns, the shear strength of the material is used. This means that the different binder recipes are assessed according to the shear strength obtained. The connection between P-wave velocity and compressive strength is indirect and material dependent, i.e., depends on the composition of the soil material. A correlation between P-wave velocity and shear strength of the stabilized soil is reported in Figure \ref{Fig:05}.

\begin{figure}[H]
 \centering
    \includegraphics[width=0.45\textwidth]{Fig_05.jpg}\\
  \caption{Shear strength as a function of P-wave velocity. (a): red circles indicate identified outliers. (b): outliers removed}\label{Fig:05}
\end{figure}

In Figure \ref{Fig:05} a), all measurements carried out in the project between the P-wave velocity and shear strength are reported. In addition to variation in binder composition, the curing time of the specimens also varies. In the graph there are some obvious outliers which have been identified and removed from the correlation plot, Figure \ref{Fig:05} b). The results from the above measurements can be used to predict shear strength based on measured P-wave velocity from the \emph{in situ} measurements, e.g., with the crosshole measurements.

\subsection{Effects from adhesives on jet columns}

In order to evaluate different binders for jet column reinforcement, strength testing was performed according to EN 196-1, see Figure \ref{Fig:06} and Figure \ref{Fig:07}. The tests performed according to EN 196-1 contain binder, water and a standard sand (Anon, 2005a). In the other tests, only binders and water were used, and in some samples a mixture of clay was used. The results show, as expected, that cement without admixture of other binders gives the highest compressive strength. Then followed a mixture of 50\% CEM II and 50\% LKD. The combinations with 50\% CEM II and 50\% CKD gave a similar result to the combination of 50\% CEM II and 30\% LKD and 20\% limestone filler. 

\begin{figure}[H]
 \centering
    \includegraphics[width=0.45\textwidth]{Fig_06.jpg}\\
  \caption{Compressive strength of specimens with different binder combinations. The test is carried out by CEMEX according to the EN 196-1 standard. PresafeP = CKD}\label{Fig:06}
\end{figure}

Figure \ref{Fig:07} shows the compressive strength as a function of cement content for different curing times. The result shows the highest strength for pure CEM II followed by the mixture 60\% CEM II and 40\% CKD. The mixture of 50\% CEM II in combination with 50\% CKD showed the lowest strength.

\begin{figure}[H]
 \centering
    \includegraphics[width=0.45\textwidth]{Fig_07.jpg}\\
  \caption{Compressive strength as a function of cement content for different curing times. Tests were carried out according to EN 196-1 standard by CEMEX}\label{Fig:07}
\end{figure}

The setting time and water requirement for different binder combinations were investigated according to EN 196-3 standard, Figure \ref{Fig:08}. Water demand was the highest for binder combinations 50\% CEM II and 50\% CKD and for 50\% CEM II, 30\% LKD and 20\% limestone filler. The water requirement for 50\% CEM II and 50\% LKD was the same as for pure CEM II. The bonding time varied from 165 to 250 minutes. The admixtures with 50\% LKD and 50\% CKD, respectively, had a significantly longer setting time compared to CEM II, Figure \ref{Fig:08}.

\begin{figure}[H]
 \centering
    \includegraphics[width=0.45\textwidth]{Fig_08.jpg}\\
  \caption{Water requirement and setting time of the different binders. PresafeP = CKD. The tests carried out by CEMEX according to standard EN 196-3}\label{Fig:08}
\end{figure}

Figure \ref{Fig:09} shows compressive strength of sample stabilized with various binder recipes with different admixtures of clay during 28 days at 20\degree C. For binder recipe A, a compressive strength between 11.5 and 19.8 MPa was obtained for specimens consisting of binder and water. 

\begin{figure*}[h!]
\centering
\includegraphics[width=0.90\textwidth]{Fig_09}
\caption{Compressive strength of sample stabilized with various binder recipes with different admixtures of clay during 28 days at 20\degree C. \textbf{A}: CEM II/A-M; \textbf{B}: 70\% CEM II/A-M and 30\% LKD; \textbf{C}: 80\% CEM II/A-M and 20\% CKD; \textbf{D}: 70\% CEM II/A-M and 30\% CKD.}
\label{Fig:09}
\end{figure*}

When clay was mixed in, the strength dropped to a minimum of 4 MPa for a mixture of 58\% clay, see Figure \ref{Fig:09}. For binder recipe B, a compressive strength between 5.9 and 13.3 MPa was obtained without mixing clay. At a clay mixture of 58\%, the compressive strength dropped to ca. 3.2 MPa at its lowest. For binder recipe C, a compressive strength between 5.3 and 8.8 MPa was obtained without clay mixing. At a clay mixture of 58\%, the compressive strength dropped to 3.6 MPa. For binder recipe D, a compressive strength between 2.8 and 4.3 MPa was obtained without clay mixing, Figure \ref{Fig:09}. 

\begin{figure}[H]
 \centering
    \includegraphics[width=0.45\textwidth]{Fig_10.jpg}\\
  \caption{Compilation of the compressive strength of the various binder combinations without mixing clay.}\label{Fig:10}
\end{figure}

Figure \ref{Fig:10} shows a comparison of the compressive strength of specimens stabilized with  different recipes without any admixture of clay. 

Figure \ref{Fig:11} shows the P-wave velocity as a function of curing time for recipes A to D.

\begin{figure}[H]
 \centering
    \includegraphics[width=0.45\textwidth]{Fig_11.jpg}\\
  \caption{P-wave velocity as a function of curing time at 20\degree C. Each P-wave velocity determination is an average of measurements on five specimens.}\label{Fig:11}
\end{figure}

Figure \ref{Fig:12} shows a compilation of measured P-wave velocities and compressive strengths for the different binder recipes.

\begin{figure}[H]
 \centering
    \includegraphics[width=0.45\textwidth]{Fig_12.jpg}\\
  \caption{Compressive strength as a function of P-wave velocity for recipes A to D. Samples stored for 28 days at 20\degree C.}\label{Fig:12}
\end{figure}

The double experiments performed show good repeatability. The result shows a significantly lower heat development for the binder combination with 50\% CEM II and 50\% LKD compared to 100\% CEM II. However, both mixtures show a similar time course from mixing until the maximum temperature is obtained. The results from the performed temperature development results for the binder slurry are shown in Figure \ref{Fig:13}.

\begin{figure}[H]
 \centering
    \includegraphics[width=0.45\textwidth]{Fig_13.jpg}\\
  \caption{Heat development with two different binder combinations as a function of time}\label{Fig:13}
\end{figure}

%\textcolor{red}{\lipsum[3-4]}

\newpage
\subsection{Simulation of seismic measurements}
Simulation of results from seismic measurements were performed in the laboratory. Socket without earth has natural frequencies in the same order of magnitude as the sample itself, which gives incorrect results, Figure \ref{Fig:14}.

\begin{figure}[H]
 \centering
    \includegraphics[width=0.45\textwidth]{Fig_14.jpg}\\
  \caption{Simulation of different modes of natural frequency for pipe sleeves}\label{Fig:14}
\end{figure}

Simulated natural frequencies for the specimens without a sleeve are shown in Figure \ref{Fig:15}.

\begin{figure}[H]
 \centering
    \includegraphics[width=0.40\textwidth]{Fig_15.jpg}\\
  \caption{Calculated natural frequencies for specimens without sleeve. The test bodies have a length of 130.7 mm.}\label{Fig:15}
\end{figure}

By simulating the test bodies placed in sleeves, other natural frequencies are obtained. The simulation shows that the measured natural frequencies originate from parts of the sleeves that do not contain any material, see Figure \ref{Fig:16}.

\begin{figure}[H]
 \centering
    \includegraphics[width=0.45\textwidth]{Fig_16.jpg}\\
  \caption{Simulated natural frequencies for sample in sleeve. Sample body 130.7 and sleeve 170 mm.}\label{Fig:16}
\end{figure}

\newpage
\section{Discussion}
The laboratory testing showed a high repeatability between trials. The spread in compressive strength/shear strength increased slightly with increased strength, which is expected as minor inaccuracies in the test specimens have a greater effect. 

The requirements for the K/C columns in situ are a minimum shear strength of 150 kPa after 28 days. To ensure this, the requirement for the laboratory packed samples is set to at least 200 kPa. Several of the mixtures that contained both cement, quicklime and LKD met the requirements of at least 200 kPa after 28 days of storage. At a storage temperature of 7\degree, final strength is reached at 80 days (560 degree days) at the earliest. 

For the jet injection, the in situ requirement is a compressive strength of 2 MPa after 28 days. For laboratory-made samples, the requirement is set to at least 3 MPa. Binder recipes A and B passed the requirement of at least 3 MPa after 28 days even with a mixture of 58\% clay. For binder recipe C, the mean compressive strength was 2.77 MPa with a clay mix of 58\%.

\newpage
\section{Conclusions}
The three different simplex experiments were performed equivalent to 80, 100 and 120 kg of binder per $m^3$ of clay. The initial tests on binder mixing and sample curing were carried out in accordance with the current Swedish standard on determination of strength in soil cemented by binders, according to SS-EN 196-1 \cite{SIS2016M} and SS-EN 196-3 \cite{SIS2016N}. 

\vspace{10pt}

\noindent
\textbf{Author contributions}: Both authors wrote text of the manuscript, performed data analysis, literature review, discussion and interpretation of the results. Dr. Per Lindh supervised the project. Both authors are responsible for the content of this manuscript, read and accepted the paper and approved its submission.

\vspace{10pt}

\noindent
\textbf{Conflict of interest}: The authors declare that the research was conducted in the absence of any commercial or financial relationships that could be construed as a potential conflict of interest

\begin{acknowledgement}
We thank the anonymous reviewers for their valuable comments. 
\end{acknowledgement}

\newpage
%\bibliographystyle{IEEEtran}
\addcontentsline{toc}{section}{References}
\bibliographystyle{plain}
\bibliography{References_NE.bib}
%\nocite*{}

\end{document}